\chapter{Background and Preliminaries}
\label{chap:background}

\section{Quantum Networks and Repeaters}
\label{sec:bg:networks}

\subsection{Bell Pairs and Entanglement Distribution}
\label{sec:bg:bellpairs}

A Bell pair is an entangled two-qubit state shared between two nodes. The
canonical Bell state is
\[
    |\Phi^+\rangle = \frac{1}{\sqrt{2}}\bigl(|00\rangle + |11\rangle\bigr).
\]
\TODO{Describe generation via photonic links, heralded Bell-state measurement
(BSM), and the role of fidelity as the key quality metric.}

\subsection{Half-RGS Repeater Architecture}
\label{sec:bg:hrgs}

\TODO{Define the half-resource-graph-state (half-RGS) building block,
the all-photonic Bell-state analyser (ABSA), and how $N$-hop chains are
assembled by joining two half-RGSs at a central node.}

\subsection{Backbone Selection and Z Side-Effect Propagation}
\label{sec:bg:backbone}

\TODO{Explain forced backbone selection at each ABSA, the
$\kappa$ span notation $(a,b)$, and why Z measurement outcomes propagate
as a constant 2-bit/ABSA/trial classical overhead.}

\section{Error Representation}
\label{sec:bg:errors}

\subsection{The Error Vector}
\label{sec:bg:errorvec}

\TODO{Define $\mathbf{e} = (w, x, y, z)^\top$ as the Pauli-channel
component vector. Explain what each component represents and why this
representation is closed under all operations in the model.}

\subsection{Z Side-Effect Algebra}
\label{sec:bg:sideffects}

\TODO{Define parity accumulation, the XOR rule across BSM outcomes,
and how the scheduling layer propagates side-effect corrections.}

\subsection{Decoherence Model}
\label{sec:bg:decoherence}

While idle, an edge's error vector relaxes toward the maximally mixed
state:
\[
    \mathbf{e}(t) = \tfrac{1}{4}\mathbf{1} +
    e^{-\gamma t}\bigl(\mathbf{e}(t_0) - \tfrac{1}{4}\mathbf{1}\bigr).
\]
\TODO{Discuss the memory coherence time $T_2 = 1/\gamma$, why corrections
are applied locally at every node rather than as a global lump, and the
three canonical timing scenarios (raw, baseline, optimistic).}

\section{Entanglement Purification}
\label{sec:bg:purification}

\subsection{Purification Circuits: YY, ZX, and XZ}
\label{sec:bg:circuits}

\TODO{Describe the three two-to-one purification protocols and which
error type each is designed to detect.}

\subsection{Success Probabilities and Output States}
\label{sec:bg:purmath}

The success probability and purified output for each circuit~\cite{bencha2025integrating} are:
\begin{align}
    P_{ZX}(\mathbf{e}_1,\mathbf{e}_2)
        &= (w_1+x_1)(w_2+x_2)+(z_1+y_1)(z_2+y_2), \label{eq:pzx}\\
    P_{XZ}(\mathbf{e}_1,\mathbf{e}_2)
        &= (w_1+z_1)(w_2+z_2)+(x_1+y_1)(x_2+y_2), \label{eq:pxz}\\
    P_{YY}(\mathbf{e}_1,\mathbf{e}_2)
        &= (w_1+y_1)(w_2+y_2)+(x_1+z_1)(x_2+z_2). \label{eq:pyy}
\end{align}
\TODO{Write out the corresponding output vectors and discuss the
legality constraint (both inputs must share the same span $\kappa$).}

\subsection{Heralded versus Optimistic Purification}
\label{sec:bg:optimistic}

\TODO{Explain the optimistic (blind) execution model: generation and
purification are pipelined without waiting for heralding; classical
outcomes are compared in a single round at the end. Show how the
position of \texttt{Herald} nodes in the schedule DAG determines
whether a branch is heralded or optimistic.}